\section{Methodology} 
We propose a two-stage method to learn a good bilingual and multilingual language-image representation model. In the first stage, we follow the work of \newcite{carlsson2022cross} to use Teacher Learning to learn a multilingual text encoder from the CLIP text encoder. In this step, no image is needed in training and only language parallel data is used. In the second stage, we use text-image pairs to further fine-tune the model from contrastive learning.  Our overall training procedure is summarized in Figure ~\ref{fig:framework}.

%\subsection{Model Architecture}
%As illustrated in Fig.~\ref{fig:framework}, our AltCLIP model has three major components. Firstly, an CLIP text encoder, an CLIP image encoder, and a student text encoder. The CLIP text encoder and CLIP image encoder are adopt from the open available CLIP models. The student text encoder is a language model trained on both English and Chinese data. 

\begin{table*}[htb]
\small
\centering
\begin{tabular}{clccccc} 
\hline
\multirow{2}{*}{Language} & \multicolumn{1}{c}{\multirow{2}{*}{Method}} & \multicolumn{5}{c}{Dataset}                                         \\ 
\cline{3-7}
                           & \multicolumn{1}{c}{}                         & ImageNet & ImageNet Sketch~ & ImageNet-A & ImageNet-R & ImageNetV2~  \\ 
\hline
\multirow{5}{*}{\rotatebox{90}{English}}   & CLIP                                         & \textbf{75.5}     & \textbf{59.6}             & \textbf{70.6}       & \textbf{87.9}       & \textbf{69.9}         \\
                           & M-CLIP                                       & 52.3     & 44.2             & 59.1       & 81.7       & 47.3         \\
                           %& M-CLIP$_B$                                      & 30       & 22.2             & 21.3       & 48.3       & 26.7         \\
                           & CN-CLIP                                      & 32.5     & 30.7             & 29.8       & 57.0         & 28.6         \\
                           & AltCLIP$_{T}$                                     & 74.7     & 59.2             & 70.4       & 87.9       & 68.8         \\
                           & AltCLIP                                      & 74.5     & 58.7             & 69.5       & 87.2       & 68.2         \\ 
\hline
\multirow{5}{*}{\rotatebox{90}{Chinese}}   & CLIP                                         & 1.9      & 1.7              & 4.7          & 4.4        & 1.8          \\
                           & M-CLIP                                       & 43.0       & 36.3             & 51.3       & 68.3       & 39.5         \\
                           %& M-CLIP$_B$                                      & 23.1     & 17.2             & 16.4       & 38.4       & 20.1         \\
                           & CN-CLIP                                      & 53.6     &47.5           & 42.8       & 78.1       & 47.8         \\
                           & AltCLIP$_{T}$                                     & 58.2     & 46.9             & \textbf{62.7}       & 82.1         & 53.3         \\
                           & AltCLIP                                      & \textbf{59.6}     &\textbf{48.4}             & 61.5       & \textbf{82.5}       &\textbf {54.0}         \\
\hline
\end{tabular}
\caption{Comparison Results of our proposed model and baseline models on image classification benchmarks, i.e. ImageNet and its variants. AltCLIP$_{T}$ and AltCLIP denotes our model after Teacher Learning Stage and after Contrastive Learning Stage, respectively. All image encoders used in these models are Vit-L for fair comparison. The metric reported is zero-shot classification accuracy.}
\label{tab:imagenet}
\end{table*}

\subsection{Teacher Learning Stage}
In this stage, we perform Teacher Learning ~\cite{hinton2015distilling} on text encoders. We use the text encoder from CLIP ~\cite{radford2021learning} as the teacher text encoder, and the XLM-R ~\cite{conneau2020unsupervised} model pretrained on multilingual data as the student encoder. A fully-connected layer is added to transform the output of the XLM-R model into the same output dimension as the teacher encoder. We use parallel text data in both English and Chinese to distill the knowledge of text-image alignment. 

Given parallel text input $(sent_1, sent_2)$, the teacher text encoder generates the learning target from input $sent_1$, which is the embedding of the \ttt{[TOS]} token, denoted by $x^{t}_{tos}$. The student text encoder generates embedding $x^{s}_{cls}$ from input $sent_2$. We minimize Mean Squared Error (MSE) between $x^{t}_{tos}$ and $x^{s}_{cls}$. 
After such training, the student text encoder can keep most of its multilingual capability and obtain text-image alignment capability in both languages. Note that the teacher encoder is only used at training time. At inference time, only the student encoder is used as the text encoder.

To show that our method is extensible at including more languages, we build a multilingual version that supports nine different languages: English(En), Chinese(Zh), Spanish(Es), French(Fr), Russian(Ru), Arabic(Ar), Japanese(Ja), Korean(Ko), and Italian(It). For the multilingual version, we align more languages with English, with the same concept and architecture as in the bilingual version.

%So, we can expand the language of the CLIP model by simply altering the text encoder into the student text encoder. 


\subsection{Contrastive Learning Stage}
This stage of training aims to further improve text-image alignment by contrastive learning on multilingual text-image pairs. As illustrated in Figure \ref{fig:framework}, here we use the image encoder from CLIP which is based on Vision Transformer (ViT) ~\cite{dosovitskiy2020image} as our image encoder, and use the student text encoder learned from the Teacher Learning Stage as our text encoder. 

We use Contrastive Loss ~\cite{hadsell2006dimensionality} between the output projection of the image encoder and text encoder, as done similarly in previous work ~\cite{radford2021learning}. We follow LiT~\cite{zhai2022lit} to freeze the image encoder at training time and only update the parameters in the text encoder. We observe that this stage of training further improves the  model's performance on various evaluation benchmarks, as presented in Section \ref{sec5}.

%\paragraph{Text Encoder}
%We use the original CLIP text encoder as our teacher text encoder. The teacher text encoder generates the learning target, the \ttt{[TOS]} token embedding $x_{tos}$. 

%The student text encoder is a multi-lingual language model that aims to align multi-language inputs to the image encoder. Here, we do not train a multi-lingual CLIP from scratch. Pretraining a non-English CLIP model from scratch is not often applicable~\cite{} due to the imbalance distribution of languages. So, we propose a cross-lingual training schema to learn the cross-lingual text-image alignment ability. After such training, the student model can keep most of its multi-lingual ability and get extra text-image alignment ability. So, we can expand the language of the CLIP model by simply altering the text encoder into the student text encoder. 

%As we know, the quality and quantity of English data are superior to other languages in the world and the English version of CLIP has superior performance than the ones in other languages. 
%Consequently, the text encoder of English CLIP achieves the best alignment ability with the image encoder. To leverage this advantage of English CLIP to expand the input languages of CLIP, we use the text encoder from the CLIP model as the teacher model. 
%By doing so, we can turn a English CLIP into a bi-lingual CLIP with fewer resources. For a given sentence $x$, the CLIP text encoder encode it into a sentence embedding as follow,
%\begin{equation}
%    T_{clip}(x) = x_{tos}\,.
%\end{equation}
%Here, $T_{clip}$ denotes the CLIP text encoder, and $x_{tos}$ represents the embedding of \ttt{[TOS]} token .
%\paragraph{Student Text Encoder}
%For a given sentence $x$, the student text encoder encode it into a sentence embedding as follow,
%\begin{equation}
%    T_{student}(x)=x_{cls}\,.
%\end{equation}
%Here, $T_{student}$ denotes the student text encoder, and $x_{cls}$ represents the embedding of \ttt{[CLS]} token.
%\paragraph{CLIP Image Encoder}
%The image encoder from CLIP model is used as our image encoder. In cross-lingual training schema, the student text encoder will be trained to align multi-language inputs to the CLIP image encoder directly. For a given image $z$, the CLIP image encoder encode it into a image embedding as follow,
%\begin{equation}
%    I_{clip}(z) = z_{cls}\,,
%\end{equation}
%where $I_{clip}$ is the CLIP image encoder, and $z_{cls}$ is the embedding of \ttt{[CLS]} token.

%\subsection{Learning Schema}
%Our training process consists of two stages. The Parallel Teacher Learning (PKD) use the parallel data to distillate the text-image alignment ability from the CLIP model. The Cross-lingual Contrastive Learning (CCL) learn the cross-ligngual text-image alignment in a contrasting manner. 

%\paragraph{Parallel Teacher Learning}
%The Parallel Teacher Learning (PKD) aims to get a multi-lingual text encoder by cross-lingual distillation. The teacher model is the CLIP text encoder which only takes English inputs. The student model  can handle both English and Chinese inputs. Inspired by previous work~\cite{}, we use the student model to learn the cross-lingual text-image alignment by Teacher Learning. We freeze the parameters of teacher model as most study dose. For convelient of training, we use a fully-connected layer to transform the output of student  model into the same output dimension of the teacher model.

%However,  previous methods can not fully take advantage of the text-image alignment of the teacher model and often get unsatisfied performance. To alleviate this problem, we propose the paralleled Teacher Learning. we can get a student model with well text-image alignment and cross-lingual ability.  We formulate the learning procedure as follow, 
%\begin{equation}
%    \ell_{pkd} = \sum_{\{x^{en},x^{cn}\}}\{mse( x^{en}_{cls}, x^{en}_{tos}) + mse(x^{cn}_{cls}, x^{en}_{tos})\}\,.
%\end{equation}
%Here, the $\ell_{pkd}$ is the parallel Teacher Learning loss, $\{x^{en},x^{cn}\}$ is from the parallel data $D_{pkd}$, $x^{en}_{tos}$ is the output of the CLIP text encoder (teacher), $x^{en}_{cls}$ and $x^{cn}_{cls}$ are both the outputs of the student text encoder.
%\paragraph{Cross-lingual contrastive Learning}
%The Cross-lingual Contrastive Learning (CCL) aims to further improve text-image alignment by contrastive learning on multi-lingual text-image pairs. We follow LiT~\cite{zhai2022lit} to freeze the Image tower of CLIP and fine-tuning the parameters in the text encoder. We observe that this further improve the model's performance on zero-shot tasks, as we present in Section \ref{}. The learning procedure is formulated as follow,
%\begin{equation}
%\begin{aligned}
%      \ell_{ccl} = &-\log \frac{\exp{(x\cdot z_+^T/\tau)}} {\sum_{i=1}^{K}\exp{(x \cdot z_i^T/\tau)}}  \\
%    &-\log \frac{\exp{(z \cdot x_+^T/\tau)}} {\sum_{i=1}^{K}\exp{(z \cdot x_i^T/\tau)}} 
%\end{aligned}
%\end{equation}
%Here, $\ell_{ccl}$ is the loss of contrastive learning, $]{x,z]}\in D_{ccl}$, $D_{ccl}$ is the dataset for contrastive learning, $K$ is the batch size. 